\documentclass[11pt,a4paper]{article}
\usepackage[margin=22mm,headheight=15pt]{geometry}
\usepackage{fontspec}
\usepackage{polyglossia}
\setdefaultlanguage{russian}
\setotherlanguage{english}
\setmainfont{Arial}
\setsansfont{Arial}
\setmonofont{Menlo}
\newfontfamily\cyrillicfont{Arial}
\newfontfamily\cyrillicfonttt{Menlo}
\usepackage{unicode-math}
\setmathfont{STIX Two Math}
\usepackage{xcolor,booktabs,longtable,array,enumitem,graphicx}
\usepackage[most]{tcolorbox}
\usepackage{fancyhdr}
\usepackage[colorlinks=true,linkcolor=blue,urlcolor=blue]{hyperref}
\definecolor{navy}{HTML}{102044}
\definecolor{blue}{HTML}{1168E8}
\definecolor{pale}{HTML}{F3F8FF}
\definecolor{yellow}{HTML}{FFDF55}
\setlength{\parindent}{0pt}
\setlength{\parskip}{5pt}
\setlist[itemize]{leftmargin=*,itemsep=2pt,topsep=4pt}
\setlist[enumerate]{leftmargin=*,itemsep=2pt,topsep=4pt}
\pagestyle{fancy}
\fancyhf{}
\fancyhead[L]{\small\color{navy} Введение в машинное обучение · Пара 2}
\fancyhead[R]{\small\color{blue} Конспект преподавателя}
\fancyfoot[C]{\small\thepage}
\renewcommand{\headrulewidth}{0.4pt}
\newtcolorbox{saybox}[1]{colback=pale,colframe=blue,boxrule=0.6pt,arc=2mm,title={#1},fonttitle=\bfseries,breakable}
\newtcolorbox{warnbox}[1]{colback=yellow!20,colframe=yellow!70!black,boxrule=0.6pt,arc=2mm,title={#1},fonttitle=\bfseries,breakable}
\newcommand{\term}[1]{\textbf{\color{navy}#1}}
\newcommand{\timing}[1]{\hfill{\small\color{blue}#1}}
\newcolumntype{P}[1]{>{\raggedright\arraybackslash}p{#1}}

\begin{document}
\begin{titlepage}
\vspace*{25mm}
{\color{blue}\bfseries\Large ДГТУ · Введение в машинное обучение\par}
\vspace{10mm}
{\color{navy}\bfseries\fontsize{30}{35}\selectfont Пара 2\\[3mm] Линейная регрессия\par}
\vspace{8mm}
{\Large Конспект для преподавателя к презентации Reveal.js\par}
\vspace{18mm}
\colorbox{yellow}{\parbox{0.86\textwidth}{\vspace{4mm}\large\textbf{Главный вопрос пары:} откуда берутся веса, если квадрат ошибки можно минимизировать точно, и какое число после этого ещё не является качеством?\vspace{4mm}}}
\vfill
{\large 90 минут · 27 слайдов · русскоязычная онлайн-лекция\par}
\vspace{5mm}
{\small Версия: октябрь 2026. Оценка качества --- по учебнику Яндекса, статьи 1.3 и 3.1--3.3. Устройство МНК --- в контуре лекций ВШЭ по линейным моделям и главы 3 ESL. Текст лекции эти источники не копирует.\par}
\end{titlepage}

\tableofcontents
\newpage

\section{Педагогическая цель и маршрут}

После занятия студент должен без подсказки: (1) записать прогноз как \(Xw\) со столбцом единиц; (2) отличить линейность по весам от линейности по сырому признаку; (3) вывести формулы наклона и сдвига для одного признака и записать нормальные уравнения; (4) объяснить, что остаток ортогонален столбцам, а при копиях столбцов единственен прогноз, не вес; (5) связать квадрат с гауссовым шумом, модуль с лапласовым и назвать четыре условия Гаусса--Маркова; (6) выбрать метрику из MSE, RMSE, MAE, \(R^{2}\), MAPE, sMAPE, RMSLE, pinball и Huber и назвать её ловушку; (7) посчитать эти числа на учебном примере и не назвать train-\(R^{2}\) качеством; (8) решить уравнение только на train и узнать утечку в матрице \(X\).

Это не лекция про градиентный спуск, KNN, полиномы и регуляризацию. Если спрашивают, как искать минимум, когда систему не решают, ответ один: неделя 4 даст шаг, неделя 5 применит его к этой же MSE.

{\small\begin{tabular}{@{}P{18mm}P{22mm}P{113mm}@{}}
\toprule
Время & Слайды & Смысл и действие преподавателя\\
\midrule
0--6 мин & 1--2 & Первый алгоритм с точным ответом. Качество --- не арифметика на тех же точках.\\
6--18 мин & 3--6 & Модель, столбец единиц, линейность по \(w\), сумма квадратов как функция весов.\\
18--40 мин & 7--9 & Вывод прямой на доске, нормальные уравнения, три точки.\\
40--52 мин & 10--11 & Проекция в \(\mathbb R^{n}\), ранг, обусловленность, SVD вместо обратной матрицы.\\
52--62 мин & 12--14 & Гаусс и Лаплас, Гаусс--Марков, чтение коэффициента. Слайды 13--14 сжимаются первыми.\\
62--80 мин & 15--19 & Loss и метрика, абсолютные и относительные метрики, квантиль, Хьюбер, логарифм ответа.\\
80--90 мин & 20--27 & Train-таблица, split, утечка, вопрос, собеседование, практика.\\
\bottomrule
\end{tabular}}

Если обсуждение затягивается, теорему Гаусса--Маркова оставьте четырьмя именами условий, а слайд про \(\exp(x^{\top}w)\) --- одной оговоркой. Не сокращайте вывод наклона, нормальные уравнения, три точки и запрет выбора по test.

\section{Гальтон и имя метода}

\timing{6--15 мин}

Слово \term{регрессия} в этом смысле ввёл Фрэнсис Гальтон. В 1885--1886 годах он сопоставил рост родителей и взрослых детей. Женский рост приводил к мужской шкале множителем около \(1{,}08\) и усреднял двух родителей: \emph{mid-parent}. Статья: Galton, \emph{Regression towards mediocrity in hereditary stature}, Journal of the Anthropological Institute, 15 (1886).

Крайние родители дают детей ближе к середине популяции. Очень высокие отцы имеют высоких сыновей, но в среднем ниже себя. Очень низкие --- низких сыновей, но в среднем выше себя. Гальтон назвал это возвратом к посредственности, \emph{regression towards mediocrity}.

На график он наносил не отдельные семьи, а средний рост сыновей в группе с данным ростом родителей. Эти средние легли почти на прямую, более пологую, чем \(45^{\circ}\). Наклон у него был около \(2/3\): отклонение родителей переходит сыну не целиком.

Современная запись при одинаковом разбросе \(x\) и \(y\):
\[
\mathbb E[y\mid x]=\mu+\beta(x-\mu),\qquad \beta=\rho.
\]
Пока \(|\rho|<1\), предсказанное отклонение меньше отклонения \(x\). Это и есть регрессия к среднему. Если дисперсии разные,
\[
\beta=\rho\,\frac{\sigma_y}{\sigma_x},
\]
и наклон бывает больше единицы. Бытовой «откат» требует \(0<\beta<1\).

Это не сила, которая тянет конкретного человека к норме. Шум, который сделал родителя крайним, детям не передаётся, поэтому группа, отобранная как крайняя, в следующем измерении ближе к середине. Тот же артефакт без биологии: худшие по первой контрольной на второй в среднем ближе к общей середине, даже если вмешательства не было. Сдвиг нельзя приписывать лечению.

\begin{warnbox}{Чего не говорить}
Метод наименьших квадратов Гальтон не вывел и Пирсон не вывел. Его опубликовал Лежандр в 1805 году; Гаусс применял его раньше. Пирсон в 1896 году записал коэффициент корреляции как формальный объект. Имя \emph{regression} принадлежит картинке Гальтона. Прямая, которую мы ищем сегодня, --- уже задача МНК для условных средних, а не биологический закон.
\end{warnbox}

\section{Модель}

\timing{6--18 мин}

Прогноз на одном объекте:
\[
\hat y = w_0 + w_1 x_1 + \cdots + w_d x_d.
\]
В векторной записи к признакам дописывают единицу, \(x=(1,x_1,\ldots,x_d)^{\top}\), и тогда \(\hat y=w^{\top}x\). На матрице из \(n\) объектов \(\hat y=Xw\). Первый столбец \(X\) состоит из единиц и отвечает за сдвиг. Без него прямая обязана проходить через начало координат: это содержательный запрет, его не включают молча. \texttt{sklearn.linear\_model.LinearRegression} добавляет столбец сам при \texttt{fit\_intercept=True}, это значение по умолчанию.

\term{Линейность здесь --- по весам.} Запись \(w_0+w_1 x+w_2 x^{2}\) остаётся линейной моделью: столбцы равны \(1\), \(x\), \(x^{2}\), веса входят в первой степени, функция \(Q(w)\) по-прежнему квадратична и выпукла. Запись \(w_0+e^{w_1}x\) линейной в этом смысле не является: нормальные уравнения для неё не написать. Полиномиальные столбцы и штраф за большие веса --- неделя 6. Сегодня достаточно развилки.

Функция, которую минимизируют на train:
\[
Q(w)=\|y-Xw\|^{2}=\sum_{i=1}^{n}(y_i-x_i^{\top}w)^{2}.
\]
Множитель \(1/n\) или \(1/2\) не меняет аргминимум. Он меняет только шкалу числа, которое называют loss. Ошибки \(1\) и \(3\) входят как \(1\) и \(9\).

\begin{saybox}{Что спросить}
Где в \(X\) живёт сдвиг? Останется ли модель линейной по весам, если вторым столбцом положить квадрат площади?
\end{saybox}

\section{Вывод для одного признака}

\timing{18--26 мин}

Пишите на доске, слайд только фиксирует ответ. Для \(\hat y=a+bx\)
\[
Q(a,b)=\sum_{i=1}^{n}(y_i-a-bx_i)^{2}.
\]
Производная по сдвигу:
\[
\frac{\partial Q}{\partial a}=-2\sum_{i=1}^{n}(y_i-a-bx_i).
\]
Ноль даёт \(\sum(y_i-a-bx_i)=0\), то есть \(a=\bar y-b\bar x\). Прямая проходит через точку \((\bar x,\bar y)\).

Производная по наклону:
\[
\frac{\partial Q}{\partial b}=-2\sum_{i=1}^{n}x_i(y_i-a-bx_i).
\]
Подставьте \(a=\bar y-b\bar x\):
\[
\sum_{i=1}^{n}x_i\bigl((y_i-\bar y)-b(x_i-\bar x)\bigr)=0.
\]
Перенос \(b\) и тождество \(\sum x_i(x_i-\bar x)=\sum(x_i-\bar x)^{2}\) (потому что \(\sum(x_i-\bar x)=0\)) дают
\[
b=\frac{\sum_{i=1}^{n}(x_i-\bar x)(y_i-\bar y)}{\sum_{i=1}^{n}(x_i-\bar x)^{2}},\qquad a=\bar y-b\bar x.
\]
Наклон --- выборочная сонаправленность, делённая на выборочный разброс \(x\). Делитель \(n\) или \(n-1\) сокращается, если он один и тот же в числителе и знаменателе.

Частный случай без признака, один столбец единиц: наклону не из чего взяться, и \(a=\bar y\). Поэтому для MSE лучшая константа --- среднее train. Для модуля производная в нуле ломается. Сдвиг константы на малый шаг уменьшает сумму модулей с одной стороны и увеличивает с другой; баланс --- когда слева и справа объектов поровну. Это медиана. Замкнутой формулы наклона для модуля нет.

\begin{warnbox}{Не смешать с первой лекцией}
На первой паре среднее было простым baseline для цены. Сегодня оно --- решение нормальных уравнений в модели из одного столбца. Какая константа лучшая, зависит от критерия. Для MAE честнее медиана. Среднее остаётся законным baseline, если в таблице так и написано, а не сказано «оптимум модуля».
\end{warnbox}

\section{Нормальные уравнения}

\timing{26--34 мин}

Остаток \(r=y-Xw\). Дифференцирование \(Q(w)=r^{\top}r\) по вектору весов даёт
\[
\nabla Q(w)=-2X^{\top}(y-Xw).
\]
Критическая точка:
\[
X^{\top}(y-Xw)=0\qquad\Leftrightarrow\qquad X^{\top}X\,w=X^{\top}y.
\]
Это и есть нормальные уравнения. Явная запись \(w=(X^{\top}X)^{-1}X^{\top}y\) требует полного столбцового ранга. На слайде её нет специально: обращать матрицу --- плохой численный рецепт, см. раздел про ранг.

Гессиан равен \(2X^{\top}X\). Для любого шага \(d\)
\[
d^{\top}(2X^{\top}X)d=2\|Xd\|^{2}\ge 0.
\]
Матрица положительно полуопределена, седла нет, критическая точка --- глобальный минимум. Если ранг неполный, минимумов много, но вектор прогнозов \(Xw\) у них один: это проекция \(y\) на пространство столбцов.

\section{Три точки}

Признак \(1,2,3\), ответы \(1,2,2\). Тогда \(\bar x=2\), \(\bar y=5/3\). Отклонения \(x\): \(-1,0,1\). Отклонения \(y\): \(-2/3,\ 1/3,\ 1/3\).
\[
b=\frac{(-1)(-2/3)+(1)(1/3)}{1+1}=\frac{1}{2},\qquad a=\frac{5}{3}-1=\frac{2}{3}.
\]
Прогнозы: \(7/6\), \(5/3\), \(13/6\). Остатки: \(-1/6\), \(+1/3\), \(-1/6\). Сумма остатков равна нулю --- следствие столбца единиц, а не удачного примера. \(\mathrm{RSS}=1/6\).

Среднее \(5/3\) как отдельная модель:
\[
\mathrm{SST}=\Bigl(-\frac{2}{3}\Bigr)^{2}+\Bigl(\frac{1}{3}\Bigr)^{2}+\Bigl(\frac{1}{3}\Bigr)^{2}=\frac{2}{3}.
\]
На этих же трёх точках, с делителем \(n=3\):

{\small\begin{tabular}{@{}lrr@{}}
\toprule
& Только среднее & Прямая МНК\\
\midrule
MSE & \(2/9\) & \(1/18\)\\
MAE & \(4/9\) & \(2/9\)\\
\(R^{2}\) & \(0\) & \(3/4\)\\
MAPE & \(1/3\) & \(5/36\)\\
\bottomrule
\end{tabular}}

\(R^{2}=1-(1/6)/(2/3)=3/4\). Прямая обязана выиграть у среднего по MSE на train: столбец единиц уже был в более узкой модели, пространство столбцов расширили. По MAE выигрыш в этом примере есть, но МНК его не гарантирует: веса искали не модулем. Все четыре строки --- арифметика обучения. Качеством их не называют.

\section{Геометрия}

\timing{40--46 мин}

Условие \(X^{\top}(y-X\hat w)=0\) означает: остаток ортогонален каждому столбцу \(X\). Картинка на плоскости \((x,y)\) --- другой объект, чем эта формула. Формула живёт в \(\mathbb R^{n}\), где \(y\) --- один вектор, а столбцы \(X\) натягивают подпространство размерности не выше числа столбцов. МНК кладёт прогноз в это подпространство так, чтобы перпендикуляр из \(y\) был кратчайшим.

Если столбец единиц есть, остаток ортогонален и ему, поэтому сумма остатков нулевая. Теорема Пифагора для центрированного ответа, а не для сырого \(y\):
\[
\|y-\bar y\mathbf 1\|^{2}=\|X\hat w-\bar y\mathbf 1\|^{2}+\|y-X\hat w\|^{2}.
\]
\(R^{2}\) --- доля первого слагаемого справа в левой части. Это доля объяснённой суммы квадратов отклонений от среднего. Не доля верных ответов, не вероятность и не «точность 75\%».

\begin{warnbox}{Какое среднее стоит в знаменателе}
На train при столбце единиц знаменатель --- сумма квадратов отклонений от среднего train, и для константной модели \(R^{2}=0\). Функция \texttt{sklearn.metrics.r2\_score} по умолчанию берёт среднее того \(y\), который передали. Если передали \(y_{\mathrm{test}}\), знаменатель живёт на test. Статья 3.2 учебника Яндекса советует константу-ориентир для MSE считать по train, чтобы не подглядывать. Это разные процедуры. На паре скажите о расхождении один раз.
\end{warnbox}

\section{Ранг, коллинеарность, численное решение}

\timing{46--52 мин}

Если два столбца совпали, \(X^{\top}X\) необратима. Семейство \(w\), дающих один и тот же \(Xw\), бесконечно: к решению можно прибавить вектор из ядра \(X\). Прогнозы на объектах, которые лежат в линейной оболочке обучающих строк, от выбора представителя не зависят. Рассказ «этот вес --- рубли за метр» зависит и поэтому запрещён.

Почти зависимые столбцы формально оставляют единственного \(w\), но делают его шумным. При сферических ошибках и полном ранге
\[
\mathrm{Var}(\hat w)=\sigma^{2}(X^{\top}X)^{-1}.
\]
Малое собственное число \(X^{\top}X\) раздувает дисперсию соответствующих комбинаций весов. Точечный прогноз \(X\hat w\) при этом может почти не сдвинуться: неустойчив рассказ про отдельный коэффициент, а не обязательно линия. Площадь и число комнат в одном уравнении --- бытовой пример. Коэффициент «за метр» не произносят отдельно от «за комнату».

Численно обращать \(X^{\top}X\) не следует. Число обусловленности \(X^{\top}X\) примерно квадрат числа обусловленности \(X\), и погрешность входа усиливается дважды. Взрослые реализации решают задачу наименьших квадратов через QR или SVD. \texttt{LinearRegression} в scikit-learn идёт этим путём и среди равных по остатку решений берёт веса минимальной нормы. Регуляризация, которая лечит почти-копии добавкой \(\lambda I\), --- неделя 6. Сегодня дыру называют.

\section{Шум, правдоподобие, Гаусс--Марков}

\timing{52--62 мин}

\subsection{Две модели ошибки}

Пусть \(y=Xw+\varepsilon\) и координаты \(\varepsilon\) независимы, \(\mathcal N(0,\sigma^{2})\). Плотность выборки пропорциональна
\[
\exp\bigl(-\|y-Xw\|^{2}/(2\sigma^{2})\bigr).
\]
Максимум правдоподобия по \(w\) совпадает с минимумом квадратов. Оценка \(\sigma^{2}\) по максимуму правдоподобия равна \(\mathrm{RSS}/n\). Несмещённая оценка при полном ранге и \(p\) столбцах --- \(\mathrm{RSS}/(n-p)\). На паре хватает одной фразы: делитель зависит от соглашения, аргминимум по \(w\) от него не зависит. Это тот же жест, что с \(n\) и \(n-1\) в дисперсии.

У распределения Лапласа плотность пропорциональна \(\exp(-|y_i-x_i^{\top}w|/b)\). Максимум правдоподобия совпадает с минимумом модуля. Для одной константы оптимум --- медиана. В общем случае замкнутой линейной системы нет; задачу решают как линейную программу. Квадрат тянет прямую к далёкой точке, потому что её вклад растёт как квадрат расстояния. Модуль растёт линейно и смотрит скорее на то, по какую сторону остатка лежит большинство точек.

\subsection{Что именно утверждает теорема}

При четырёх условиях оценка МНК --- лучшая линейная несмещённая оценка \(w\) (\emph{best linear unbiased estimator}). «Лучшая» значит: минимальная дисперсия среди оценок, которые линейны по \(y\) и несмещены. Не среди всех мыслимых оценок.

\begin{enumerate}
\item Спецификация: \(\mathbb E[y\mid X]=Xw\). Среднее и правда лежит в выбранных столбцах.
\item Экзогенность: \(\mathbb E[\varepsilon\mid X]=0\). Признаки не несут остаточной связи с ошибкой.
\item Сферичность: \(\mathrm{Var}(\varepsilon\mid X)=\sigma^{2}I\). Дисперсия одна и та же, ошибки не коррелируют друг с другом.
\item Полный столбцовый ранг. Иначе отдельный вес не идентифицирован.
\end{enumerate}

Гауссов шум в этот список не входит. Он нужен, чтобы МНК совпал с максимумом правдоподобия и чтобы точные интервалы Стьюдента имели названное распределение. Для BLUE хватает четырёх пунктов выше.

Гетероскедастичность ломает сферичность. При сохранённой экзогенности веса часто остаются несмещёнными и перестают быть эффективными; обычные стандартные ошибки, посчитанные как \(\sigma^{2}(X^{\top}X)^{-1}\), врут. Эндогенность ломает второй пункт: пропущенный признак, который сидит и в столбце, и в ошибке, смещает веса. Это уже не «шум в ответе», а другая задача. Теорема не выдаёт причинность.

\begin{saybox}{Если осталось сорок секунд}
Назовите четыре условия и одну поломку: веер дисперсии не обязан сместить прямую, но лишает МНК звания лучшей линейной несмещённой оценки. Пропущенный общий фактор смещает.
\end{saybox}

\section{Коэффициент}

Фраза, которую можно произнести: при прочих равных столбцах модели плюс одна единица признака \(j\) сдвигает прогноз на \(w_j\).

Фраза, которую произнести нельзя: признак причиняет ответ; вес устойчив, хотя столбцы ходят вместе; после стандартизации \(w_j\) всё ещё измерен в рублях за метр. Стандартизованный столбец меняет единицу: вес относится к одному стандартному отклонению. Дамми района --- разница прогноза относительно того уровня, которого в столбцах нет, а не «вклад района вообще».

\section{Метрики}

\timing{62--80 мин}

\subsection{Две роли}

\term{Функция потерь} даёт число, по которому на train ищут \(w\). \term{Метрика} отвечает, насколько прогноз полезен для решения, и её считают на строках, не участвовавших в этом поиске. Учить MSE и сообщать MAE --- нормальная пара. Совпадение формулы роли не сливает.

\subsection{Абсолютные}

\[
\mathrm{MSE}=\frac{1}{m}\sum_{j=1}^{m}(y_j-\hat y_j)^{2},\qquad
\mathrm{RMSE}=\sqrt{\mathrm{MSE}},\qquad
\mathrm{MAE}=\frac{1}{m}\sum_{j=1}^{m}|y_j-\hat y_j|.
\]
RMSE возвращает единицы \(y\), но крупные ошибки по-прежнему входят через квадрат. Лучшая константа MSE --- среднее, лучшая константа MAE --- медиана.

\[
R^{2}=1-\frac{\sum(y_j-\hat y_j)^{2}}{\sum(y_j-\bar y)^{2}}.
\]
Единица --- совпадение с ответами. Ноль --- не лучше константы среднего в выбранной записи. Отрицательное значение на новых данных --- хуже этой константы.

Скорректированный коэффициент на той же выборке,
\[
R^{2}_{\mathrm{adj}}=1-(1-R^{2})\frac{n-1}{n-p},
\]
уменьшает долю за лишние столбцы. Это штраф на train, а не замена validation. Сравнивать им модели с разным преобразованием \(y\) нельзя: знаменатель живёт в другой шкале.

\subsection{Относительные и их смещение}

\[
\mathrm{MAPE}=\frac{1}{m}\sum_{j=1}^{m}\frac{|y_j-\hat y_j|}{|y_j|}.
\]
На одной точке ошибка «вдвое выше» и «вдвое ниже» даёт один и тот же MAPE. По набору разных \(y\) это уже не так. Для константы \(c\) субградиент суммы \(|y_i-c|/|y_i|\) уравновешен, когда сумма весов \(1/|y|\) слева от \(c\) равна сумме справа. Оптимум --- взвешенная медиана с весами \(1/|y|\). Большой ответ весит меньше, константа тянется к меньшим \(y\). При \(y=0\) метрика не определена. Отрицательные ответы требуют отдельной договорённости, которой у обычного MAPE нет.

sMAPE делит модуль на \(|y|+|\hat y|\) (часто ещё с множителем \(2\)). Слово «симметричная» описывает формулу, не оптимум. При фиксированном модуле ошибки больший прогноз увеличивает знаменатель и уменьшает штраф, поэтому оптимум снова съезжает вниз.

\[
\mathrm{RMSLE}=\sqrt{\frac{1}{m}\sum_{j=1}^{m}\bigl(\log(1+y_j)-\log(1+\hat y_j)\bigr)^{2}}.
\]
Метрика сравнивает порядки. При ответе \(100\) прогноз \(50\) даёт больший логарифмический зазор, чем прогноз \(150\): при том же модуле занижение весит больше завышения. Отрицательный \(y\) в \(\log(1+y)\) не подставить.

\subsection{Квантиль и Хьюбер}

Pinball-потеря уровня \(\tau\in(0,1)\):
\[
\rho_{\tau}(r)=\begin{cases}\tau r,& r\ge 0,\\ (\tau-1)r,& r<0.\end{cases}
\]
При \(\tau=1/2\) это \(|r|/2\), и константа-оптимум --- медиана. В общем случае минимум ожидаемого pinball по константе --- квантиль уровня \(\tau\): субградиент меняет знак, когда доля точек слева достигает \(\tau\). Если недопрогноз дороже перепрогноза, берут \(\tau>1/2\). Это уже функция потерь под асимметричное решение, не ещё один синоним MAE.

Хьюбер с порогом \(\delta>0\): квадрат при \(|r|\le\delta\) и линейный рост дальше,
\[
H_{\delta}(r)=\begin{cases}\tfrac12 r^{2},& |r|\le\delta,\\ \delta(|r|-\delta/2),& |r|>\delta.\end{cases}
\]
В нуле функция гладкая, в отличие от модуля, и хвост не получает квадратичного голоса. Это loss для обучения. В отчёт бизнесу её не кладут. \(\delta\) выбирают по validation; на test порог не крутят.

\subsection{Логарифм ответа}

Если учат \(\log y=x^{\top}w+\varepsilon\) при \(\varepsilon\sim\mathcal N(0,\sigma^{2})\) и \(y>0\), то условное распределение \(y\) логнормально. Его медиана равна \(\exp(x^{\top}w)\), а среднее ---
\[
\mathbb E[y\mid x]=\exp(x^{\top}w+\sigma^{2}/2).
\]
Поправка \(\sigma^{2}/2\) принадлежит этой модели шума. Её не ставят универсальным коэффициентом «чтобы перевести логарифм в рубли». Неравенство Йенсена объясняет направление без нормальности: \(\exp\) выпукла, поэтому среднее \(\exp(\varepsilon)\) больше единицы, если \(\mathbb E\varepsilon=0\) и \(\varepsilon\) не постоянно. Метрика в логарифме и метрика в рублях отвечают на разные вопросы. В таблице эксперимента единица измерения написана явно, иначе строка с \(\log y\) побеждает строку в рублях видом числа.

\section{Где решают уравнение}

\timing{80--90 мин}

\begin{itemize}
\item \term{Train} --- столбцы, веса, среднее константы, \(\mu\) и \(\sigma\) любых преобразований.
\item \term{Validation} --- какой набор столбцов оставить, какую метрику отчёта выбрать, какой \(\delta\) у Хьюбера.
\item \term{Test} --- один отчёт уже выбранного решения.
\end{itemize}

Случайный \texttt{train\_test\_split} честен, только если строки независимы и одинаково устроены. Время не перемешивают: будущая цена не проверяется перемешанным будущим. Одна сущность не стоит по обе стороны split. \texttt{random\_state} фиксируют. \(k\)-fold на этой паре не открывают.

Утечка здесь --- конкретная поломка матрицы, а не общий лозунг. Три дыры практики:

\begin{enumerate}
\item \term{Прокси ответа.} Цена за метр, умноженная на площадь, восстанавливает цену. Нормальные уравнения найдут почти нулевой остаток и пустые веса.
\item \term{Статистика до split.} Среднее для заполнения пропуска, \(\mu\) и \(\sigma\) масштаба, словарь категорий оценивают на train и затем применяют к остальным частям:
\[
z=\frac{x-\mu_{\mathrm{train}}}{\sigma_{\mathrm{train}}}.
\]
На train значения \(0,0,10\), на test значение \(10\). \(\mu_{\mathrm{train}}=10/3\), среднее всех четырёх точек равно \(5\). Отложенный объект уже вошёл в преобразование.
\item \term{Одна сущность дважды.} Кадры одного видео, сделки одного дома, один человек в train и test. Модель запоминает сущность. Тот же механизм у случайного split по времени.
\end{enumerate}

Высокий score --- повод проверить список, не доказательство дыры и не доказательство качества.

Таблица эксперимента --- одна строка на законченный протокол. На validation из отдельного, не досочного примера истинные цены \(6\) и \(14\), константа train равна \(10\) и даёт MAE \(4\); прямая даёт прогнозы \(7\) и \(13\), MAE \(1\). Двух точек мало, чтобы объявить модель годной. Третью строку, где масштаб посчитан по всей матрице, рядом не ставят: у неё другой протокол.

\begin{saybox}{Три финальных вопроса}
1. Какая константа минимизирует MSE, и из какого частного случая нормальных уравнений она берётся?\\
2. Почему совпавшие столбцы убивают рассказ про вес, но не обязательно прогноз?\\
3. Чем MAPE отличается от MAE, кроме деления на \(y\)?
\end{saybox}

Ответы: среднее, потому что в модели из столбца единиц критическая точка --- \(\bar y_{\mathrm{train}}\); прогнозы --- проекция на пространство столбцов, и прибавление вектора из ядра \(X\) её не меняет; MAPE не определена при нуле и тянет константу к меньшим ответам как взвешенную медиану с весами \(1/|y|\).

Вопрос собеседования: \(R^{2}\) на train равен \(0{,}98\), хороша ли модель. Сильный ответ отказывается. Train-\(R^{2}\) растёт от добавления столбцов; число ничего не говорит ни о новых объектах, ни о том, совпадает ли метрика с решением, ни о копиях столбцов и прокси ответа внутри \(X\). Слабый ответ: «это 98\% точности».

\section{Шпаргалка: частые ошибки}

{\small\begin{longtable}{@{}P{58mm}P{99mm}@{}}
\toprule
Неверное утверждение & Что ответить\\
\midrule
\endhead
«Квадрат признака выводит модель из линейных» & По весам она остаётся линейной. Ломает линейность вес внутри функции.\\
«Делим на \(n\) или на \(n-1\) --- получим другие веса» & Для аргминимума квадрата множитель не важен. Делитель важен для оценки \(\sigma^{2}\).\\
«Всегда пишем \(w=(X^{\top}X)^{-1}X^{\top}y\)» & Это формула при полном ранге. Считают через QR или SVD.\\
«Коллинеарность портит прогнозы, значит и веса» & Часто наоборот: веса разъезжаются, проекция остаётся устойчивой.\\
«\(R^{2}=0{,}75\) --- точность 75\%» & Доля объяснённых квадратов против выбранной константы.\\
«MAPE симметрична, раз в названии нет квадрата» & На разных \(y\) оптимум смещён к меньшим ответам. При нуле метрики нет.\\
«sMAPE чинит смещение MAPE» & Знаменатель растёт вместе с прогнозом, оптимум снова съезжает вниз.\\
«\(\exp\) от логарифмического прогноза --- ожидаемая цена» & Для нормального шума в логарифме это медиана. Среднее больше на \(\exp(\sigma^{2}/2)\).\\
«Гаусс--Марков требует нормальных ошибок» & Нормальность нужна для ММП. Для BLUE нужны спецификация, экзогенность, сферичность и ранг.\\
«Высокий \(R^{2}\) доказывает отсутствие утечки» & Он совместим и с прокси ответа, и с одной сущностью по обе стороны split.\\
\bottomrule
\end{longtable}}

\section{Откуда взяты идеи}

Оценка качества и запрет смешивать выборки собраны по программе курса и по статьям учебника Яндекса, без переноса их текста:

\begin{itemize}
\item \href{https://contest.yandex.ru/tracks/ml/introduction/machine-learning}{1.3. Машинное обучение}: выбор метрики под цель, запрет смешивать train и test.
\item \href{https://contest.yandex.ru/tracks/ml/model-evaluation/introduction}{3.1. Введение в оценку качества}: метрика как критерий сравнения; важно, на каких объектах она посчитана.
\item \href{https://contest.yandex.ru/tracks/ml/model-evaluation/classification-and-regression-metrics}{3.2. Метрики}: функция потерь и метрика; MSE, RMSE, \(R^{2}\) относительно константы; MAE при хвостах.
\item \href{https://contest.yandex.ru/tracks/ml/model-evaluation/cross-validation}{3.3. Кросс-валидация}: hold-out, shuffle, прокси ответа, статистики до split, группы и время.
\end{itemize}

Устройство МНК --- стандартный контур курса линейных моделей ВШЭ (Е.~А.~Соколов) и главы 3 \emph{The Elements of Statistical Learning}: нормальные уравнения, проекция остатка, выпуклость квадрата, различие замкнутой формы и итеративного поиска. Теорема Гаусса--Маркова названа в своей классической границе, без обещания причинности. На этой паре не выводим \(k\)-fold, градиентный шаг, Ridge и метрики классификации.

\section{Памятка по слайдам}

{\footnotesize\begin{longtable}{@{}P{8mm}P{48mm}P{97mm}@{}}
\toprule
№ & Слайд & Переход\\
\midrule
\endhead
1 & Линейная регрессия & KNN, спуск и Ridge не обещать.\\
2 & Откуда берутся веса & Пример 1, 2, 2 при \(x=1,2,3\) назван заранее.\\
3 & Взвешенная сумма & Единица уже внутри \(x\), если есть сдвиг.\\
4 & Столбец единиц & Без него прямая проходит через ноль.\\
5 & Линейность по весам & \(x^{2}\) можно, \(e^{w}\) нельзя.\\
6 & Сумма квадратов & \(1/n\) не двигает минимум.\\
7 & Формулы \(a\) и \(b\) & Вывод на доске. Константа --- среднее.\\
8 & Нормальные уравнения & Градиент, гессиан, единственность при полном ранге.\\
9 & Три точки & \(a=2/3\), \(b=1/2\), сумма остатков нулевая.\\
10 & Ортогональность & Пифагор для центрированного \(y\). \(R^{2}\) --- не точность.\\
11 & Ранг & Прогноз единственен, вес нет. SVD, не обратная матрица.\\
12 & Гаусс и Лаплас & ММП. Делитель \(\sigma^{2}\) --- одна фраза.\\
13 & Гаусс--Марков & Четыре условия. Нормальность не входит. Можно сжать.\\
14 & Чтение \(w_j\) & При прочих равных в модели. Не причина.\\
15 & Loss и метрика & Учить MSE, сообщать MAE.\\
16 & MSE, RMSE, MAE, \(R^{2}\) & Adjusted --- штраф на train. Оговорить sklearn.\\
17 & MAPE, sMAPE, RMSLE & Каждая смещает оптимум по-своему.\\
18 & Pinball и Хьюбер & \(\tau=1/2\) --- медиана. \(\delta\) только по validation.\\
19 & Логарифм ответа & \(\exp(x^{\top}w)\) --- медиана. Можно свернуть.\\
20 & Таблица на трёх точках & Train. По MSE прямая обязана выиграть.\\
21 & Три роли & Время, группа, запрет правки по test.\\
22 & Три дыры & Прокси, статистика, сущность. \(\mu=10/3\) против \(5\).\\
23 & Строка эксперимента & 6 и 14 --- другой пример. Третью строку не сравнивать.\\
24 & Какое утверждение верно & Ответ Б.\\
25 & \(R^{2}=0{,}98\) & Дождаться отказа и двух слоёв аргумента.\\
26 & Практика & Руки, sklearn, четыре метрики, утечка, таблица.\\
27 & Итог & Три вопроса из раздела про split.\\
\bottomrule
\end{longtable}}

\end{document}
